%----------------------------------------------------------------------------------------
%	PACKAGES AND OTHER DOCUMENT CONFIGURATIONS
%----------------------------------------------------------------------------------------
\DocumentMetadata{
pdfversion=2.0,  
lang=en-US,   
pdfstandard={ua-2, a-4f},
tagging = on, 
tagging-setup={math/setup={mathml-SE}} ,
}
\documentclass[12pt]{article}
\usepackage{amsmath}
\usepackage[extreme]{savetrees}
\usepackage[utf8]{inputenc}
\usepackage{newpxtext}
\usepackage{hyperref}
\usepackage{multicol}
\usepackage{physics}
\usepackage[customcolors]{hf-tikz}

%%%%%%%%% === Document Configuration === %%%%%%%%%%%%%%

\hypersetup{
pdftitle={MATH 340 -- Handout 1 Spring 2026},
pdfauthor={Ignacio Rojas},
pdfkeywords={first handout},%
}

\tikzset{
  style green/.style={
    set fill color=green!50!lime!60,
    set border color=white,
  },
  style cyan/.style={
    set fill color=cyan!90!blue!60,
    set border color=white,
  },
  style orange/.style={
    set fill color=orange!80!red!60,
    set border color=white,
  },
  row/.style={
    above left offset={-0.15,0.31},
    below right offset={0.15,-0.125},
    #1
  },
  col/.style={
    above left offset={-0.1,0.3},
    below right offset={0.15,-0.15},
    #1
  }
}

\newcommand{\twobyone}[2]{\begin{pmatrix} %% 2 x 1 matrix
  #1 \\ #2 \end{pmatrix}}
\newcommand{\twobytwo}[4]{\begin{pmatrix} %% 2 x 2 matrix
    #1 & #2 \\ #3 & #4 \end{pmatrix}}
\newcommand{\twobythree}[6]{\begin{pmatrix} %% 2 x 3 matrix
        #1 & #2 & #3\\ #4 & #5 & #6 \end{pmatrix}}
\newcommand{\threebyone}[3]{\begin{pmatrix} %% 3 x 1 matrix
  #1 \\ #2 \\ #3 \end{pmatrix}}
\newcommand{\threebytwo}[6]{\begin{pmatrix} %% 3 x 1 matrix
    #1 & #2\\ #3 & #4\\ #5&#6 \end{pmatrix}}
\newcommand{\threebythree}[9]{\begin{pmatrix} %% 3 x 3 matrix
  #1 & #2 & #3 \\ #4 & #5 & #6 \\ #7 & #8 & #9 \end{pmatrix}}

\newcommand{\To}{\Rightarrow}

%----------------------------------------------------------------------------------------
%	ARTICLE CONTENTS
%----------------------------------------------------------------------------------------

\begin{document}
\begin{multicols}{2}
\textbf{Attention to notation!}
\begin{itemize}
  \item Swap: Exchange rows.
  \begin{itemize}
    \item \(R_i=R_j,\)
    \item \(R_i\leftrightarrow R_j,\) 
    \item \(r_i\leftrightarrow r_j.\)
  \end{itemize}
  \item Scaling: Multiply rows by constants.
  \begin{itemize}
    \item \(cR_i\),
    \item \(R_i=cR_i\), 
    \item \(cr_i\).
  \end{itemize}
  \item Combination: Add multiples of rows.
  \begin{itemize}
    \item \(R_i+cR_j,\)
    \item \(R_i=R_i+cR_j,\)
    \item \(R_i\leftarrow R_i+cR_j.\)
  \end{itemize}
\end{itemize}

After applying this row operations, systems of equations keep having the same solutions. If a system is augmented, the augmentation also is affected by the same row operations.

\textbf{Remark:}
Given that systems are \emph{equivalent} after row operations, we apply them in order to get simpler equations to solve.

\textbf{Definition:}
A matrix is in \underline{row echelon form} if:
\begin{enumerate}
  \item The zero rows are at the \textbf{bottom}.
  \item The \(1^{\text{st}}\) nonzero entry of each row \textbf{is to the right} of the \(1^{\text{st}}\) nonzero entry of the previous row.
\end{enumerate}


\textbf{Example:}
Consider the matrix 

\[A=\begin{pmatrix}
  1&0&1&-2&0\\ 0&0&0&-1&3\\ 0&0&2&4&-5
\end{pmatrix}.\]

Let's see if \(A\) is in row echelon form.\par 
 
\begin{itemize}
  \item[C1] Condition 1 of the definition is met because there are no zero rows.
  \item[C2] The \(1^{\text{st}}\) nonzero entries of each row are 1, -1 and 2. The -1 in the second row is indeed to the \textbf{right} of the 1. But the 2 is to the \textbf{left} of the -1.
\end{itemize}
\vfill\null\columnbreak
Condition 2 is not met, so \(A\) is not in row echelon form. But, can we convert it? Yes! With a row operation. Notice that by applying a \textbf{swap} we obtain:
\[
\begin{pmatrix}
  1&0&1&-2&0\\ 0&0&0&-1&3\\ 0&0&2&4&-5
\end{pmatrix}\xrightarrow[]{R_2\leftrightarrow R_3}\begin{pmatrix}
  1&0&1&-2&0\\ 0&0&2&4&-5\\ 0&0&0&-1&3
\end{pmatrix}.
\]
The matrix \(A\) and the new matrix \textbf{are not the same}. However, this new matrix \textbf{is} in row echelon form.


\textbf{Practice:}
Consider the matrix 
\[
B=\begin{pmatrix}
  1&2\\1&3\\0&0\\0&1
\end{pmatrix}.\]
 Is it in row echelon form? If not, convert it using row operations!\par As a \textbf{hint}, you should use the \textbf{combination} operation. 

\textbf{Definition:}
  A matrix is in \underline{reduced} \underline{row echelon form} (\textbf{RREF}) if:
  \begin{enumerate}
    \item It is in \textbf{row echelon form}.
    \item The \(1^{\text{st}}\) nonzero entry of each row is a 1.
    \item In each \textbf{column}, the leading 1s are \textbf{only accompanied by zeros}.
  \end{enumerate}


\textbf{Example:}
  Consider the last matrix we obtained from the previous example.
  \[\begin{pmatrix}
    1&0&1&-2&0\\ 0&0&2&4&-5\\ 0&0&0&-1&3
  \end{pmatrix}\]
  Is it in \textbf{reduced row echelon form}?
  \begin{enumerate}
    \item[C1] Yes, it is in row echelon form.
    \item[C2] The \(1^{\text{st}}\) nonzero entries of each row are 1, 2 and -1. So only the first row meets this condition.
  \end{enumerate}
  We remedy this by applying two \textbf{scalings}:
  \begin{align*}
    \begin{pmatrix}
      1&0&1&-2&0\\ 0&0&2&4&-5\\ 0&0&0&-1&3
    \end{pmatrix}&\xrightarrow{\frac12 R_2}\begin{pmatrix}
      1&0&1&-2&0\\ 0&0&1&2&\frac{-5}{2}\\ 0&0&0&-1&3\\
    \end{pmatrix}\\
      &\xrightarrow{-1 R_3}\begin{pmatrix}
        1&0&1&-2&0\\ 0&0&1&2&\frac{-5}{2}\\ 0&0&0&1&-3\\
    \end{pmatrix}
  \end{align*}
  \vfill\null\columnbreak
  With this, condition 2 is now met. However, condition 3 is still not met. The 1 in the second row has a 1 above it, and the one in the third row has a -2 and a 2. To get rid of them, we apply a \textbf{combination}.
  \begin{align*}
    \begin{pmatrix}
        1&0&1&-2&0\\ 0&0&1&2&\frac{-5}{2}\\ 0&0&0&1&-3\\
    \end{pmatrix}&\xrightarrow{R_1-R_2} \begin{pmatrix}
      1&0&0&-4&\frac{5}{2}\\ 0&0&1&2&\frac{-5}{2}\\ 0&0&0&1&-3\\
  \end{pmatrix}\\
  &\xrightarrow{R_2-2R_3} \begin{pmatrix}
      1&0&0&-4&\frac{5}{2}\\ 0&0&1&0&\frac{7}{2}\\ 0&0&0&1&-3\\
\end{pmatrix}\\
&\xrightarrow{R_1+4R_3} \begin{pmatrix}
  1&0&0&0&\frac{-19}{2}\\ 0&0&1&0&\frac{7}{2}\\ 0&0&0&1&-3\\\end{pmatrix}
  \end{align*}

  Condition 3 is thus met, and this is the reduced row echelon form of \(A\).

\textbf{Remark:}
  By saying \emph{the} reduced row echelon form, we allude to uniqueness. This is because \textbf{every matrix} is row-equivalent to a \textbf{unique} matrix in \textbf{RREF}.


  \textbf{Practice:} Let's solve the following system of equations in variables \((a,b,c,d)\):
  \[
  \left\lbrace\begin{aligned}
    &3a-c=0\\
    &8a-2d=0\\
    &2b-2c-d=0
              \end{aligned}\right.
  \]

  \textbf{Practice:} Consider the system of equations in variables \((u,v,w,x,y)\) below:
  \[
  \left\lbrace\begin{aligned}
    &u+v-y=1\\
    &v+2w+x+3y=1\\
    &u-w+x+y=0
              \end{aligned}\right.
  \]
 
  \begin{enumerate}
    \item What is the coefficient matrix associated with the system? And the augmented matrix?
    \item Reduce the augmented matrix to find the solution. ?
    \item Which variables are free in the solution? How many parameters does the solution depend on? 
  \end{enumerate}

\subsection*{The Matrix Product}

The condition that guarantees that two matrices \(A,B\) can be multiplied is:
\
If \(A\) is \([m\times p]\) and \(B\) is \([p\times n]\), then the entries of the product are defined by the formula:
\
More simply:
\
The matrix \(AB\) will be of size \([m\times n]\).

\textbf{Example:}
Consider the matrices

\[
A=\begin{pmatrix}
    1&-1\\
    3&4\\
    -2&0\\
    0&4
\end{pmatrix},\ B=\begin{pmatrix}
    2&0&-2&8&-5\\
    1&-2&3&9&6
\end{pmatrix}.
\]
\vfill\null\columnbreak

Notice that:
\begin{itemize}
    \item \(A\) is \([4\times 2]\).
    \item \(B\) is \([2\times 5]\).
    \begin{itemize}
        \item It holds that \(\#\text{cols. } A = \#\text{rows }B\).
        \item We can multiply \(A\) and \(B\)!
    \end{itemize}
    \item The matrix \(AB\) will be of size \([4\times 5]\).
\end{itemize}
If we want to find the \((2,4)\) entry of \(AB\), we use row 2 of \(A\) and column \(4\) of \(B\):


\[
\left(\begin{array}{cc}
    1&-1\\
    \tikzmarkin[row=style cyan]{r2}3&4\tikzmarkend{r2}\\
    -2&0\\
    0&4
  \end{array}\right)\quad\text{y}\quad
  \left(\begin{array}{ccccc}
    2&0&-2&\tikzmarkin[col=style green]{c4}8&-5\\
    1&-2&3&9\tikzmarkend{c4}&6
  \end{array}\right).
\]

We take the dot product of these vectors:
\[(3,4)\cdot(8,9)=3\cdot 8+4\cdot 9=60\]
and therefore \((AB)_{24}=60.\)


In general, we can arrange the matrices in the following way:

\[
\begin{array}{cc|c|c|c|c|c}
    &&2&0&-2&\tikzmarkin[col = style green]{c4d2}8&-5\\
    &&1&-2&3&9&6\\
    \hline1&-1&&&&&\\
    \hline\tikzmarkin[row = style cyan]{r2d2}3&4&&&&(\ast)\tikzmarkend{r2d2}\tikzmarkend{c4d2}&\\
\hline-2&0&&&&&\\
   \hline 0&4&&&&&\\
\end{array}
\]


\textbf{Practice:}
  \begin{itemize}
    \item Find the \((1,3)\) entry of \(AB\) by highlighting the corresponding rows and columns.
    \item Perform the same procedure to find the \((3,2)\) entry of \(AB\).
  \end{itemize}

%https://tex.stackexchange.com/questions/232516/cannot-compile-matrix-multiplication-example-from-texample-net
%https://tex.stackexchange.com/questions/448843/matrix-multiplication-with-tikz
To find any entry, we locate the corresponding row and column of the matrices and calculate their dot product.



\textbf{Exercise:}
We already multiplied \(A\) on the left of \(B\). Can we multiply \(BA\)? That is, reversing the order of multiplication.


\textbf{Remark:}
In general, if \(A\) has dimensions \([m\times n]\) and \(B\) has dimensions \([n\times m]\), then both \(AB\) and \(BA\) exist.


\textbf{Example:}
  Consider the matrices 
  \[
  A=\begin{pmatrix}
    1&-1\\
    3&4\\
    -2&0\\
    0&4
\end{pmatrix},\ B=\begin{pmatrix}
    2&0&-2&8&-5\\
    1&-2&3&9&6
\end{pmatrix}.
  \]

If we want to find \(AB\), we first verify if we can multiply them. \(A\) is \([2\times 4]\) and \(B\) is \([4\times 2]\); in this case, 
\[
\text{\#rows } A=\text{\#cols. } B
\]
so we can calculate \(AB\).\par 
In the same way, if we look for \(BA\), we can also calculate it because \[
\text{\#rows } B=\text{\#cols. } A.
\]


\textbf{Practice:}
  \begin{enumerate}
    \item What size will the matrices \(AB\) and \(BA\) have?
    \item Work together with someone to find the resulting matrices from the multiplication.
  \end{enumerate}

\subsection*{Determinants}

\textbf{Notation:}
The determinant of a matrix \(A\) is denoted by \(\det(A)\) or \(|A|\).


\begin{itemize}
  \item In the \([2\times 2]\) case, if \(A=\begin{pmatrix}
    a&b\\c&d
  \end{pmatrix}\), then
  \[\det(A)=ad-bc.\]
  
  \item The \([3\times 3]\) case, \textbf{Sarrus' Rule}:\par 
  If \(A=\begin{pmatrix}
    a&b&c\\d&e&f\\g&h&i
  \end{pmatrix}\), then we calculate its determinant by taking the first two columns and augmenting \(A\) with them such that we obtain
  \[\left(\begin{array}{ccc|cc}
    a&b&c&a&b\\d&e&f&d&e\\g&h&i&g&h
  \end{array}\right)\]
  from here we add the diagonals from top to bottom and subtract the diagonals from bottom to top such that 
  \[\det(A)=aei+bfg+cdh-gec-hfa-idb.\]
\end{itemize}

\subsection*{Eigenvalues and Eigenvectors}

Eigenvectors are generators for invariant spaces of linear maps; that is, \emph{they don't change direction after the mapping.}

\textbf{Definition:} An \underline{eigenvector} of a map \(T\) is a non-zero vector \(\vec{x}\) that satisfies \(T(\vec{x})=\lambda\vec{x}\) for some scalar \(\lambda\). The constant \(\lambda\) is called the \underline{eigenvalue} associated with \(\vec{x}\).\par

To find the eigenvectors, we modify the defining equation \(T(\vec{x})=\lambda\vec{x}\) as follows:

\begin{align*}
  T(\vec{x})=\lambda\vec{x}&\iff T(\vec{x})=\lambda\operatorname{id}(\vec{x})\\
  &\iff T(\vec{x})-\lambda\operatorname{id}(\vec{x})=0\\
  &\iff \left(T-\lambda\operatorname{id}\right)(\vec{x})=0\\
  &\iff \left(A-\lambda I\right)(\vec{x})=0
\end{align*}

So we are tasked with finding the solutions \(\vec{x}\) to the matrix equation:
\[
 \left(A-\lambda I\right)(\vec{x})=0.
\]

\textbf{Remark:} For the matrix in question to have non-zero solutions to the homogeneous equation, we must guarantee that it is non-invertible. So that: 
\[
A-\lambda I\ \text{is non-invertible}\iff \det(A-\lambda I)=0.
\]

\textbf{Definition:} For a linear map \(T\) represented by a matrix \(A\), we define its \underline{characteristic polynomial} as: 
\[
f_A(x)=\det(A-xI).
\]
The zeroes, or roots, of the characteristic polynomial will be \(T\)'s eigenvalues.

\textbf{Remark:}
Observe that geometrically, we often start by identifying the eigenvectors, because invariant subspaces can be visually perceived. However, algebraically we proceed in reverse: finding the eigenvalues first and then their corresponding eigenvectors.

\textbf{Definition:} 
The \underline{algebraic multiplicity} (a.m.) of an eigenvalue \(\lambda\) is the degree of the factor \((x-\lambda)\) in the characteristic polynomial. It is also the number of times \(\lambda\) appears as an eigenvalue of the map.\par
The \underline{geometric multiplicity} (g.m.) of an eigenvalue \(\lambda\) is the number of linearly independent eigenvectors associated with \(\lambda\).

\textbf{Important:} 
It always holds that: 
\[1\leq \text{g.m.}\leq \text{a.m.}\]
In other words: every eigenvalue always has at least one eigenvector associated with it. Furthermore, there can be no more linearly independent eigenvectors than instances of the eigenvalue they are associated with.

\textbf{Example:} 
We have \(C=\twobytwo{1}{3}{0}{1}\), a shear mapping. 
\begin{itemize}
  \item \(C\) is triangular \(\To\) \(\lambda=1\) with a.m.\(=2\).
  \item It has a single eigenvector \((1,0)\), so g.m.\(=1\).
  \item Since a.m. \(\neq\) g.m., \(C\) is not diagonalizable.
\end{itemize}

%https://ece.uwaterloo.ca/~dwharder/Integer_eigenvalues/

\textbf{Example:}
  Let's consider 
  \[A=\twobytwo{2}{1}{1}{2}.\] 
  We are looking for its eigenvectors.
  \begin{enumerate}
    \item Using the trace-determinant formula we have that: 
    \[m=\frac{1}{2}(2+2),\ p=\det A=3\To \lambda=2\pm\sqrt{2^2-3},\]
    so \(\lambda_1=3\) and \(\lambda_2=1\).
    
    \item The first eigenspace is: 
    \[E_{\lambda_1}=\ker(A-3I)=\ker\twobytwo{-1}{1}{1}{-1}=\ker\twobytwo{1}{-1}{0}{0}.\]
    The equation associated with this matrix is \(x-y=0\) and its solution is: 
    \[(x,y)=(x,x)=x(1,1)\To E_{\lambda_1}=\operatorname{span}(1,1).\]
    We conclude that \(\vec{v}_1=(1,1)\) is an eigenvector of \(A\) associated with \(\lambda_1=3\).
    
    \item The second eigenspace is: 
    \[E_{\lambda_2}=\ker(A-I)=\ker\twobytwo{1}{1}{1}{1}=\ker\twobytwo{1}{1}{0}{0}.\]
    The associated equation is \(x+y=0\). Its solution is: 
    \[(x,y)=(x,-x)=x(1,-1)\To E_{\lambda_2}=\operatorname{span}(1,-1).\]
    and thus the other eigenvector is \(\vec{v}_2=(1,-1)\) associated with \(\lambda_2=1\).
  \end{enumerate}
 
  
 \textbf{Example:}
  Let's consider 
  \[A=\threebythree{0}{0}{-2}{0}{2}{2}{-2}{2}{1}.\] We are looking for its eigenvectors.
  \begin{enumerate}
    \item To find the eigenvalues, we must find the roots of the characteristic polynomial. We compute the determinant: 
    \[f(\lambda) = \det(A-\lambda I) = \det\threebythree{-\lambda}{0}{-2}{0}{2-\lambda}{2}{-2}{2}{1-\lambda}.\]
    To expand this determinant, we can use Sarrus' rule. We augment the matrix with its first two columns:
    \[\left(\begin{array}{ccc|cc}
      -\lambda & 0 & -2 & -\lambda & 0 \\
      0 & 2-\lambda & 2 & 0 & 2-\lambda \\
      -2 & 2 & 1-\lambda & -2 & 2
    \end{array}\right)\]
    Next, we sum the products of the top-to-bottom diagonals and subtract the products of the bottom-to-top diagonals:
    \begin{align*}
      f(\lambda) &= [ (-\lambda)(2-\lambda)(1-\lambda) + (0)(2)(-2) + (-2)(0)(2) ] \\
      &\quad - [ (-2)(2-\lambda)(-2) + (2)(2)(-\lambda) + (1-\lambda)(0)(0) ] \\
      &= [ -\lambda(2 - 3\lambda + \lambda^2) + 0 + 0 ] - [ 4(2-\lambda) - 4\lambda + 0 ] \\
      &= (-\lambda^3 + 3\lambda^2 - 2\lambda) - (8 - 4\lambda - 4\lambda) \\
      &= (-\lambda^3 + 3\lambda^2 - 2\lambda) - (8 - 8\lambda)\\
      &= -\lambda^3 + 3\lambda^2 + 6\lambda - 8.
    \end{align*}
    By factoring the polynomial, we obtain \(f(\lambda) = -(\lambda - 1)(\lambda - 4)(\lambda + 2)\), so our eigenvalues are \(\lambda_1=1\), \(\lambda_2=4\), and \(\lambda_3=-2\). Because all three eigenvalues are distinct, we already know that a.m.\(=1\) for each of them.
    
    \item The first eigenspace is: 
    \[E_{\lambda_1}=\ker(A-I)=\ker\threebythree{-1}{0}{-2}{0}{1}{2}{-2}{2}{0}.\]
    We find the reduced row echelon form (RREF) of this matrix:
    \[
    \threebythree{1}{0}{2}{0}{1}{2}{0}{0}{0}
    \]
    This yields the equations \(x+2z=0\) and \(y+2z=0\). Its solution is: 
    \begin{align*}
      (x,y,z)&=(-2z,-2z,z)\\
      &=z(-2,-2,1)\\
      \To E_{\lambda_1}&=\operatorname{span}(-2,-2,1).
    \end{align*}
    We conclude that \(\vec{v}_1=(-2,-2,1)\) is an eigenvector of \(A\) associated with \(\lambda_1=1\).
    
    \item The second eigenspace is: 
    \[E_{\lambda_2}=\ker(A-4I)=\ker\threebythree{-4}{0}{-2}{0}{-2}{2}{-2}{2}{-3}.\]
    Finding its RREF gives:
    \[
    \threebythree{1}{0}{1/2}{0}{1}{-1}{0}{0}{0}
    \]
    This yields the equations \(x+\frac{1}{2}z=0\) and \(y-z=0\). To avoid fractions in our vector, we can set \(z=-2x\), which then gives \(y=-2x\). Its solution is: 
    \begin{align*}
      (x,y,z)&=(x,-2x,-2x)\\
      &=x(1,-2,-2)\\
      \To E_{\lambda_2}&=\operatorname{span}(1,-2,-2).
    \end{align*}
    Thus, the second eigenvector is \(\vec{v}_2=(1,-2,-2)\) associated with \(\lambda_2=4\).

    \item The third eigenspace is: 
    \[E_{\lambda_3}=\ker(A+2I)=\ker\threebythree{2}{0}{-2}{0}{4}{2}{-2}{2}{3}.\]
    Finding its RREF gives:
    \[
    \threebythree{1}{0}{-1}{0}{1}{1/2}{0}{0}{0}
    \]
    This yields the equations \(x-z=0\) and \(y+\frac{1}{2}z=0\). To avoid fractions, we set \(z=-2y\) and thus \(x=-2y\). Its solution is: 
    \begin{align*}
      (x,y,z)&=(-2y,y,-2y)\\
      &=y(-2,1,-2)\\
      \To E_{\lambda_3}&=\operatorname{span}(-2,1,-2).
    \end{align*}
    And thus the final eigenvector is \(\vec{v}_3=(-2,1,-2)\) associated with \(\lambda_3=-2\).
  \end{enumerate}
  We see that for all three eigenvalues, a.m.\(=\)g.m.\(=1\).

 \subsection*{Generalized eigenvectors}

The geometric idea of a generalized eigenspace is that it is an entire line, plane, or subspace \emph{which doesn't pass through the origin} (\emph{it is an affine space!}) that, when acted upon by a linear map, gets pushed into a standard eigenspace.

\textbf{Definition:} If \(\vec{v}_1\) is a standard eigenvector associated with \(\lambda\), a \underline{generalized eigenvector} \(\vec{v}_2\) is a vector that gets pushed into the direction of \(\vec{v}_1\). Algebraically, we find it by solving the non-homogeneous matrix equation:
\[(A-\lambda I)\vec{v}_2 = \vec{v}_1\]

\textbf{Example:}
We have the shear mapping from before, 
\[C=\twobytwo{1}{3}{0}{1},\]
we know:
\begin{itemize}
  \item \(C\) is triangular \(\To\) \(\lambda=1\) with a.m.\(=2\).
  \item It has a single eigenvector \(\vec{v}_1=(1,0)\), so g.m.\(=1\).
  \item Since a.m. \(\neq\) g.m., \(C\) has a generalized eigenvectors.
\end{itemize}

To find the generalized eigenvector \(\vec{v}_2\), we solve the equation \((C - I)\vec{v}_2 = \vec{v}_1\):
\[\twobytwo{0}{3}{0}{0}\begin{pmatrix}x\\y\end{pmatrix} = \begin{pmatrix}1\\0\end{pmatrix}\]
This yields the equation 
\[
3y=1\To y=1/3
\]
and \(x\) is free. So that the solution to the system is
\[
(x,y)=(x,1/3)=x(1,0)+(0,1/3).
\] 
The solution includes the eigenvector itself but also the generalized eigenvector which we have not seen before.
Thus, \(\vec{v}_2=(0,1/3)\) is a generalized eigenvector for \(C\).

\textbf{Example:}
Let's consider \(A=\threebythree{0}{-1}{-1}{1}{2}{0}{0}{0}{2}\). We will find all its eigenvectors and generalized eigenvectors.
\begin{enumerate}
  \item First, we find the characteristic polynomial:
  \[f(\lambda) = \det(A-\lambda I) = \det\threebythree{-\lambda}{-1}{-1}{1}{2-\lambda}{0}{0}{0}{2-\lambda}.\]
  Expanding along the third row gives:
  \begin{align*}
      f(\lambda) &= (2-\lambda) \det\twobytwo{-\lambda}{-1}{1}{2-\lambda} \\
      &= (2-\lambda)(-\lambda(2-\lambda) - (-1)) \\
      &= (2-\lambda)(\lambda^2 - 2\lambda + 1) \\
      &= (2-\lambda)(\lambda - 1)^2.
  \end{align*}
  Our eigenvalues are \(\lambda_1=1\) (with a.m.\(=2\)) and \(\lambda_2=2\) (with a.m.\(=1\)).
  
  \item Let's find the eigenspace for \(\lambda_1=1\):
  \[E_{\lambda_1}=\ker(A-I)=\ker\threebythree{-1}{-1}{-1}{1}{1}{0}{0}{0}{1}.\]
  We find the RREF of this matrix:
  \[\threebythree{1}{1}{0}{0}{0}{1}{0}{0}{0}\]
  This yields the equations \(x+y=0\) and \(z=0\). Its solution is:
  \begin{align*}
    (x,y,z) &= (-y, y, 0) \\
    &= y(-1, 1, 0) \\
    \To E_{\lambda_1} &= \operatorname{span}(-1, 1, 0).
  \end{align*}
  We conclude that \(\vec{v}_1=(-1,1,0)\) is the only regular eigenvector for \(\lambda_1=1\). Since g.m.\(=1\) but a.m.\(=2\), we need a generalized eigenvector!
  
  \item To find the generalized eigenvector \(\vec{v}_2\) associated with \(\lambda_1=1\), we solve \((A-I)\vec{v}_2 = \vec{v}_1\). We set up the augmented matrix:
  \[\left(\begin{array}{ccc|c} -1 & -1 & -1 & -1 \\ 1 & 1 & 0 & 1 \\ 0 & 0 & 1 & 0 \end{array}\right)\]
  Finding its RREF gives:
  \[\left(\begin{array}{ccc|c} 1 & 1 & 0 & 1 \\ 0 & 0 & 1 & 0 \\ 0 & 0 & 0 & 0 \end{array}\right)\]
  This yields the equations \(x+y=1\) and \(z=0\). From here
  \[
  y=1-x
  \]
  so that our solution looks like
  \begin{align*}
    (x,y,z) &= (x, 1-x, 0) \\
   &=(-x)(-1,1,0)+(0,1,0)
  \end{align*}
  Thus, \(\vec{v}_2=(0,1,0)\) is our generalized eigenvector and we see the original eigenvector in the solution.
  
  \item Finally, we find the eigenspace for \(\lambda_2=2\):
  \[E_{\lambda_2}=\ker(A-2I)=\ker\threebythree{-2}{-1}{-1}{1}{0}{0}{0}{0}{0}.\]
  Finding its RREF gives:
  \[\threebythree{1}{0}{0}{0}{1}{1}{0}{0}{0}\]
  This yields the equations \(x=0\) and \(y+z=0\). Its solution is:
  \begin{align*}
    (x,y,z) &= (0, -z, z) \\
    &= z(0, -1, 1) \\
    \To E_{\lambda_2} &= \operatorname{span}(0, -1, 1).
  \end{align*}
  And thus the third regular eigenvector is \(\vec{v}_3=(0,-1,1)\) associated with \(\lambda_2=2\).
  \vfill\null\columnbreak
\end{enumerate}
\vfill\null
\end{multicols}
\end{document} 